\chapter{Local Perturbations in Image Space}
\label{ch:rq3}

Global perturbations mix the response associated with one token's corresponding
image region with responses caused by changes elsewhere. This chapter perturbs
one image patch and measures how far assignment changes spread beyond it.
Chapter~\ref{ch:rq2} showed assignment instability under global noise; the
present experiment tests whether changes also appear outside a bounded patch.

Finite convolutional receptive fields create overlapping local dependencies,
while attention and spatial normalization can couple distant positions. The
evaluated VQGAN contains $3\times3$ stride-2 downsampling, $3\times3$ residual
convolutions, GroupNorm, and attention at its lowest-resolution level. These
facts identify candidate mechanisms, not the cause of the measured response.

\section{Experimental Setup}
\label{sec:rq3_setup}

\subsection{Patch-Noise Protocol}
\label{subsec:rq3_patch_noise_protocol}

Let $x \in [-1,1]^{H \times W \times 3}$ be an input image and let $P$ denote
a square patch. The binary mask $M_P$ equals one inside the patch and zero
outside it. For a noise level $\sigma$, the perturbed image is
\begin{equation}
    x_{\sigma,P}
    =
    \mathrm{clip}\!\left(x + \sigma \epsilon \odot M_P, -1, 1\right),
    \qquad
    \epsilon \sim \mathcal{N}(0,I).
    \label{eq:rq3_patch_noise}
\end{equation}
Only pixels inside $P$ receive additive noise. The clipping operation keeps the image in the valid pixel range.

Both tokenizers produce a $16\times16$ token grid. The analysis uses an
$8\times8$ token-space patch, corresponding to a $128\times128$ pixel region.
Its top-left corner is sampled uniformly from valid grid locations for each of
the 50,000 ImageNet validation images. The reported levels are
$\sigma\in\{0.1,0.25,0.5\}$. The two runs use the same images and nominal
settings but different random patch locations and noise realizations. Strictly
paired model comparisons would require shared perturbations and a rerun.

The masked variant fills the selected patch with black in the internal image
range. With fill value $c=-1$, it is defined as
\begin{equation}
    x_{\mathrm{mask},P}=(1-M_P)\odot x + M_P\odot c.
    \label{eq:rq3_patch_mask}
\end{equation}

\subsection{Inside and Outside Response}
\label{subsec:rq3_inside_outside_metrics}

Let $z_0 \in \{1,\ldots,K\}^{16 \times 16}$ be the clean token grid and $z_{\sigma,P}$ be the token grid obtained from the patch-perturbed image. Let $I_P$ be the set of token positions covered by the patch and let $O_P$ be the complement. We define the inside response as the fraction of token assignments within the selected patch that change:
\begin{equation}
    R_{\mathrm{in}}(\sigma,P)
    =
    \frac{1}{|I_P|}
    \sum_{i \in I_P}
    \mathbf{1}\!\left[z_{\sigma,P,i} \neq z_{0,i}\right].
    \label{eq:rq3_inside_response}
\end{equation}
We define the outside response as the fraction of token assignments outside the selected patch that change:
\begin{equation}
    R_{\mathrm{out}}(\sigma,P)
    =
    \frac{1}{|O_P|}
    \sum_{i \in O_P}
    \mathbf{1}\!\left[z_{\sigma,P,i} \neq z_{0,i}\right].
    \label{eq:rq3_outside_response}
\end{equation}
We define the encoder leakage ratio as
\begin{equation}
    \lambda_{\mathrm{enc}}(\sigma,P)
    =
    \frac{R_{\mathrm{out}}(\sigma,P)}
         {R_{\mathrm{in}}(\sigma,P) + \varepsilon}
    \label{eq:rq3_leakage_ratio}
\end{equation}
where $\varepsilon>0$ prevents division by zero. All responses are fractions of
token assignments that differ from the clean encoding. A value of
$\lambda_{\mathrm{enc}}$ near zero indicates that changes are mostly confined
to the patch, while a value near one indicates similar change rates inside and
outside it. We retain $R_{\mathrm{in}}$ and $R_{\mathrm{out}}$ alongside the
ratio because the same ratio can arise from different absolute responses.

\subsection{Distance Profile}
\label{subsec:rq3_distance_metrics}

The inside and outside responses do not show how far token changes spread. Each
token position $i$ is therefore assigned a Manhattan grid distance $d(i,P)$ to
the nearest token in $I_P$: distance zero is inside the patch, distance one is
the first ring of neighboring token cells, and larger values move farther
across the token grid. We define the distance profile as
\begin{equation}
    R_d(\sigma,P)
    =
    \frac{1}{|\{i : d(i,P)=d\}|}
    \sum_{i : d(i,P)=d}
    \mathbf{1}\!\left[z_{\sigma,P,i} \neq z_{0,i}\right].
    \label{eq:rq3_distance_profile}
\end{equation}
The results report the full profile and a small set of bands: inside tokens ($d=0$), boundary-neighbor tokens ($d=1$), near tokens ($d \in \{1,2\}$), and far-field tokens ($d \geq 6$).

\section{Experimental Results}
\label{sec:rq3_results}

\paragraph{Qualitative patch perturbations.}
\label{subsec:rq3_qualitative}

Figure~\ref{fig:rq3_qual_combined} shows selected examples for the two tokenizers. Each panel contains the clean image, the patch-noisy inputs, the corresponding reconstructions, and binary token-response maps. The black rectangle marks the perturbed patch in pixel space and the corresponding token block in the response map. Highlighted cells outside the rectangle are token-assignment changes outside the patch.

\begin{figure}[tbp]
    \centering
    \textbf{(a) VQGAN}\par\smallskip
    \includegraphics[width=\textwidth,height=0.35\textheight,keepaspectratio]{ch06_fig_patch_noise_qualitative_vqgan.png}
    \begin{minipage}{0.96\textwidth}
        \footnotesize\centering
        Sample ILSVRC2012\_val\_00000293:
        $\sigma=0.1$ ($R_{\mathrm{in}}=0.766$, $R_{\mathrm{out}}=0.297$, $\lambda_{\mathrm{enc}}=0.388$);
        $\sigma=0.25$ ($0.969$, $0.630$, $0.651$);
        $\sigma=0.5$ ($0.969$, $0.724$, $0.747$).
    \end{minipage}

    \medskip
    \textbf{(b) LlamaGen}\par\smallskip
    \includegraphics[width=\textwidth,height=0.35\textheight,keepaspectratio]{ch06_fig_patch_noise_qualitative_llamagen.png}
    \begin{minipage}{0.96\textwidth}
        \footnotesize\centering
        Sample ILSVRC2012\_val\_00000293:
        $\sigma=0.1$ ($R_{\mathrm{in}}=0.875$, $R_{\mathrm{out}}=0.323$, $\lambda_{\mathrm{enc}}=0.369$);
        $\sigma=0.25$ ($0.938$, $0.526$, $0.561$);
        $\sigma=0.5$ ($0.984$, $0.656$, $0.667$).
    \end{minipage}
    \caption{Selected local patch-noise examples for (a) VQGAN and (b) LlamaGen using the same input ID. The binary response maps use dark cells for unchanged assignments (0) and yellow cells for changed assignments (1); the marked rectangle identifies the perturbed patch. Metrics below each panel are sample-specific and ordered as $R_{\mathrm{in}}$, $R_{\mathrm{out}}$, and $\lambda_{\mathrm{enc}}$.}
    \label{fig:rq3_qual_combined}
\end{figure}
\paragraph{Local response and leakage.}
\label{subsec:rq3_inside_outside_results}

Table~\ref{tab:rq3_inside_outside_summary} reports the mean inside response
$R_{\mathrm{in}}$, outside response $R_{\mathrm{out}}$, encoder leakage ratio
$\lambda_{\mathrm{enc}}$, and far-field response. A high inside response is
expected because the corresponding image pixels are directly perturbed. At
$\sigma=0.1$, VQGAN changes 78.8\% of inside assignments and LlamaGen changes
82.2\%. At $\sigma=0.5$, the inside response is almost saturated for both.

Token assignments also change outside the patch at the lowest tested level:
$R_{\mathrm{out}}$ is 29.9\% for VQGAN and 30.2\% for LlamaGen at
$\sigma=0.1$. It rises to 44.1\% and 42.8\%, respectively, at $\sigma=0.25$,
then to 52.1\% and 50.5\% at $\sigma=0.5$. At the largest tested value,
outside assignments change about half as often as inside assignments. A
lower-noise sweep would be needed to estimate when $R_{\mathrm{out}}$
approaches zero.

\begin{table}[t]
    \centering
    \small
    \begin{tabular}{lrrrrr}
        \textbf{Tokenizer} & $\boldsymbol{\sigma}$ & $\boldsymbol{R_{\mathrm{in}}}$ & $\boldsymbol{R_{\mathrm{out}}}$ & $\boldsymbol{\lambda_{\mathrm{enc}}}$ & \textbf{Far field} \\ \hline
        VQGAN & 0.10 & 0.788 & 0.299 & 0.375 & 0.174 \\
        VQGAN & 0.25 & 0.952 & 0.441 & 0.462 & 0.265 \\
        VQGAN & 0.50 & 0.988 & 0.521 & 0.527 & 0.333 \\ \hline
        LlamaGen & 0.10 & 0.822 & 0.302 & 0.361 & 0.178 \\
        LlamaGen & 0.25 & 0.960 & 0.428 & 0.444 & 0.260 \\
        LlamaGen & 0.50 & 0.991 & 0.505 & 0.509 & 0.328 \\
    \end{tabular}
    \caption{Inside and outside token response under local patch-noise perturbations. The far-field column reports the mean token flip rate for positions at distance at least six token cells from the perturbed patch.}
    \label{tab:rq3_inside_outside_summary}
\end{table}
\begin{figure}[t]
    \centering
    \includegraphics[width=0.92\textwidth]{ch06_fig_inside_outside_response.png}
    \caption{Mean token response inside and outside the perturbed patch, with the corresponding leakage ratio. Outside-token changes increase monotonically with noise level for both tokenizers, and the outside/inside ratio also increases.}
    \label{fig:rq3_inside_outside_response}
\end{figure}
\paragraph{Distance from the perturbed patch.}
\label{subsec:rq3_distance_results}

Figure~\ref{fig:rq3_distance_profile_flip} plots the response to the image
perturbation as a function of distance from the patch. It declines sharply over
the first few token cells but remains nonzero farther away. Table~\ref{tab:rq3_distance_band_summary}
summarizes the curves using descriptive distance bands. The chosen far-field
band, $d \geq 6$, separates distant positions from the immediate rings; it is
not a measured boundary between local and global mechanisms. In this band,
17.4\% of VQGAN assignments and 17.8\% of LlamaGen assignments change at
$\sigma=0.1$. At $\sigma=0.5$, the fraction is roughly one third.

\begin{table}[t]
    \centering
    \small
    \begin{tabular}{lrrrr}
        \textbf{Tokenizer} & $\boldsymbol{\sigma}$ & \textbf{Inside} & \textbf{Boundary} & \textbf{Far field} \\ \hline
        VQGAN & 0.10 & 0.788 & 0.597 & 0.174 \\
        VQGAN & 0.25 & 0.952 & 0.827 & 0.265 \\
        VQGAN & 0.50 & 0.988 & 0.917 & 0.333 \\ \hline
        LlamaGen & 0.10 & 0.822 & 0.613 & 0.178 \\
        LlamaGen & 0.25 & 0.960 & 0.832 & 0.260 \\
        LlamaGen & 0.50 & 0.991 & 0.922 & 0.328 \\
    \end{tabular}
    \caption{Distance-band summary of token response. Inside denotes $d=0$, boundary denotes $d=1$, and far field denotes positions with Manhattan distance $d \geq 6$ from the perturbed patch.}
    \label{tab:rq3_distance_band_summary}
\end{table}
\begin{figure}[t]
    \centering
    \includegraphics[width=0.86\textwidth]{ch06_fig_distance_profile_flip.png}
    \caption{Token flip probability as a function of distance from the perturbed patch. The response is strongest inside and near the patch, but remains measurable far from it.}
    \label{fig:rq3_distance_profile_flip}
\end{figure}
\paragraph{Masked patch intervention.}
\label{subsec:rq3_masked_results}

The masked intervention defined in Equation~\ref{eq:rq3_patch_mask} replaces the
same local patch with black. Table~\ref{tab:rq3_masked_summary} shows that this
nearly saturates the inside response for both tokenizers, while the response to
the image perturbation outside the patch is 56.3\% for VQGAN and 49.5\% for
LlamaGen. Thus, a change confined to one quarter of the image changes token
assignments across roughly half of the remaining grid.

\begin{table}[t]
    \centering
    \small
    \begin{tabular}{lrrr}
        \textbf{Tokenizer} & $\boldsymbol{R_{\mathrm{in}}}$ & $\boldsymbol{R_{\mathrm{out}}}$ & $\boldsymbol{\lambda_{\mathrm{enc}}}$ \\ \hline
        VQGAN & 0.998 & 0.563 & 0.564 \\
        LlamaGen & 1.000 & 0.495 & 0.495 \\
    \end{tabular}
    \caption{Inside and outside token response for the masked-patch intervention. The mask is a stronger local corruption than Gaussian patch noise and produces a correspondingly high outside response.}
    \label{tab:rq3_masked_summary}
\end{table}
\section{Discussion}
\label{sec:rq3_discussion}

Both tokenizers show the same measured pattern: token changes are most frequent
inside and near the perturbed patch, decline over the nearby token cells, and
remain nonzero at distant positions. LlamaGen has a slightly higher inside
response at $\sigma=0.1$ and a slightly lower leakage ratio at each tested
value, but the spatial pattern is shared.

The evaluated architecture contains overlapping $3\times3$ convolutions,
GroupNorm, and an attention layer. A small feature movement can also cross a
nearest-neighbor quantization boundary and change the assignment at that
position, although pointwise quantization does not transmit changes between
positions. The measurements do not isolate the contribution of these
components. Future work can investigate attention, spatial normalization, and
receptive-field overlap through ablations and intermediate-feature analysis.

This result complements Chapter~\ref{ch:rq2}. Global noise changed many token
assignments when every image region was perturbed. The study of perturbation
locality shows how far those changes spread when the image modification is
bounded: token assignments change well beyond the selected patch. These
changes also alter the token sequences supplied to an autoregressive model.
